% !TeX root = ../GMCM-main.tex
\section{问题重述}
\subsection{问题背景}
在具身智能、自然人机交互和智能服务系统等领域，情感分析构成了不可或缺的关键支撑。在真实人机交互过程中，人的情绪并非由单一的形态信号进行表达，而是同时通过语言内容、语言韵律、面部神态以及身体动作等多个信号进行协同体现的。因此，仅依赖文本、语音或视觉中的某一模态进行情感判断，往往难以准确刻画复杂场景下的真实情感状态。多模态信息联合建模能够充分利用不同模态之间的互补性与一致性，是突破单模态的识别局限，进而提升复杂场景下对情绪理解的准确性和稳定性。
然而，在实际情景下，采集过程中更容易受到环境噪声、画面遮挡、设备故障和语音转写失败等因素影响，致使某一模态在局部时间段内缺失。若采用常规固定权重进行模态融合，可能会把缺失值误认为有效信息，进而造成预测结果性能的下降。与此同时，若仅给出情感类别或情感强度而不说明判断依据，则会降低模型在智能交互场景中的可复核性。基于上述背景，本题以CMU-MOSEI英文多模态情感数据为基础，针对进行多模态特征的提取与时序对齐、缺失性预测和可解释性情感这三方面进行研究。
\subsection{问题提出}
结合题目要求，需要解决以下问题：
\begin{enumerate}
  \item 问题一：填写需要求解的对象、条件与目标。
  \item 问题二：填写需要分析或优化的内容。
  \item 问题三：填写需要验证或推广的内容。
\end{enumerate}
